\documentclass[11pt]{article}
\usepackage{amsmath,amssymb,amsthm}
\usepackage{algorithm}
\usepackage{algpseudocode}
\usepackage{fullpage}
\usepackage{hyperref}
\usepackage{enumitem}
\newtheorem{theorem}{Theorem}
\newtheorem{lemma}{Lemma}
\newtheorem{assumption}{Assumption}
\newtheorem{corollary}{Corollary}
\newcommand{\E}{\mathbb{E}}
\newcommand{\R}{\mathbb{R}}
\newcommand{\norm}[1]{\left\|#1\right\|}
\newcommand{\inner}[1]{\left\langle#1\right\rangle}
\newcommand{\grad}{\nabla}
\begin{document}

\section*{Algorithm Name}
\textbf{Differentially Private Stochastic Gradient Descent with Gaussian Noise (DP‑SGD)}  

The algorithm performs stochastic gradient descent while adding independent Gaussian noise to each (mini‑batch) gradient in order to satisfy $(\varepsilon,\delta)$‑differential privacy.

\section*{Mathematical Formulation}
Let $f:\R^{d}\rightarrow\R$ be the (possibly non‑convex) loss function, and let $\mathcal{B}_{t}\subset\{1,\dots,n\}$ denote the random mini‑batch sampled at iteration $t$ of size $b$.  
Define the (clipped) stochastic gradient
\[
g_{t}\;:=\;\frac{1}{b}\sum_{i\in\mathcal{B}_{t}}\operatorname{clip}\!\bigl(\grad f_{i}(x_{t}),\,C\bigr)\;+\;z_{t},
\]
where 
\begin{itemize}[noitemsep,topsep=0pt]
\item $\operatorname{clip}(v,C)=v\cdot \min\!\bigl\{1,\;C/\|v\|\bigr\}$ bounds the $\ell_{2}$‑norm of each per‑example gradient by $C>0$,
\item $z_{t}\sim\mathcal{N}\!\bigl(0,\sigma^{2}I_{d}\bigr)$ is isotropic Gaussian noise,
\item $\eta_{t}>0$ is the learning rate.
\end{itemize}
The update rule is
\begin{equation}
x_{t+1}=x_{t}-\eta_{t}g_{t}.
\tag{1}
\end{equation}
The noise scale $\sigma$ is chosen according to a privacy accountant (e.g., moments accountant) to guarantee a total privacy budget $(\varepsilon,\delta)$ after $T$ iterations.

\section*{Pseudocode}
\begin{algorithm}[H]
\caption{DP‑SGD with Gaussian Noise}
\begin{algorithmic}[1]
\Require Data $\{(x_i,y_i)\}_{i=1}^{n}$, loss $f$, clipping norm $C$, noise scale $\sigma$, learning rates $\{\eta_{t}\}_{t=0}^{T-1}$, batch size $b$, total iterations $T$.
\State Initialize $w_{0}\in\R^{d}$ (e.g., $w_{0}=0$).
\For{$t=0$ \textbf{to} $T-1$}
    \State Sample a mini‑batch $\mathcal{B}_{t}$ of size $b$ uniformly without replacement.
    \State Compute per‑example gradients $g_{i}^{t}=\grad f_{i}(w_{t})$, $i\in\mathcal{B}_{t}$.
    \State Clip: $\tilde g_{i}^{t}= \operatorname{clip}(g_{i}^{t},C)$.
    \State Average: $\bar g_{t}= \frac{1}{b}\sum_{i\in\mathcal{B}_{t}}\tilde g_{i}^{t}$.
    \State Sample noise $z_{t}\sim\mathcal{N}(0,\sigma^{2}I_{d})$.
    \State Form private gradient $g_{t}= \bar g_{t}+z_{t}$.
    \State Update: $w_{t+1}= w_{t}-\eta_{t} g_{t}$.
\EndFor
\State \Return $w_{T}$ (or averaged iterate $\frac{1}{T}\sum_{t=0}^{T-1}w_{t}$).
\end{algorithmic}
\end{algorithm}

\section*{Dry Run Example}
Consider the scalar quadratic $f(w)=(w-3)^{2}$, whose exact gradient is $\grad f(w)=2(w-3)$.  
Take $w_{0}=0$, learning rate $\eta=0.1$, clipping norm $C=10$ (so clipping never activates), and deterministic “noise” $z_{t}=0$ for illustration (the privacy accountant would set $\sigma$ later).  

\begin{center}
\begin{tabular}{c|c|c|c}
$t$ & $g_{t}=\grad f(w_{t})$ & $w_{t+1}=w_{t}-\eta g_{t}$ & $f(w_{t+1})$\\\hline
0 & $2(0-3)=-6$ & $0-0.1(-6)=0.6$ & $(0.6-3)^{2}=5.76$\\
1 & $2(0.6-3)=-4.8$ & $0.6-0.1(-4.8)=1.08$ & $(1.08-3)^{2}=3.6864$\\
2 & $2(1.08-3)=-3.84$ & $1.08-0.1(-3.84)=1.464$ & $(1.464-3)^{2}=2.365\\
\end{tabular}
\end{center}
After three iterations the objective has decreased from $9$ to approximately $2.37$, illustrating the descent despite the added (here zero) noise.

\section*{Assumptions}
\begin{assumption}[Smoothness]\label{as:smooth}
$f$ is $L$‑smooth: for all $x,y\in\R^{d}$,
\[
f(y)\le f(x)+\inner{\grad f(x)}{y-x}+\frac{L}{2}\|y-x\|^{2}.
\]
\end{assumption}
\begin{assumption}[Bounded Gradient Moments]\label{as:bounded}
For every data point $i$ and any $x$, the (clipped) stochastic gradient satisfies
\[
\|\operatorname{clip}(\grad f_{i}(x),C)\|\le C,
\qquad
\E\!\left[\|\operatorname{clip}(\grad f_{i}(x),C)-\grad f(x)\|^{2}\right]\le \zeta^{2},
\]
where the expectation is over the random sampling of a single example.
\end{assumption}
\begin{assumption}[Unbiasedness of Mini‑batch]\label{as:unbiased}
Conditioned on $x_{t}$, the averaged clipped gradient $\bar g_{t}$ satisfies
\[
\E[\bar g_{t}\mid x_{t}]=\grad f(x_{t}).
\]
\end{assumption}
\begin{assumption}[Noise Independence]\label{as:noise}
The Gaussian noises $\{z_{t}\}_{t\ge0}$ are independent of each other and of the data sampling, with
\[
z_{t}\sim\mathcal{N}\bigl(0,\sigma^{2}I_{d}\bigr),\qquad \E[z_{t}]=0,\;\E[\|z_{t}\|^{2}]=d\sigma^{2}.
\]
\end{assumption}
\begin{assumption}[Learning Rate]\label{as:lr}
The step size is constant $\eta_{t}\equiv\eta$ with $0<\eta\le \frac{1}{L}$.
\end{assumption}

\section*{Main Convergence Theorem}
\begin{theorem}[Expected Gradient Norm Decrease]\label{thm:main}
Under Assumptions \ref{as:smooth}–\ref{as:lr}, after $T$ iterations of DP‑SGD the following bound holds:
\[
\frac{1}{T}\sum_{t=0}^{T-1}\E\!\bigl[\|\grad f(x_{t})\|^{2}\bigr]
\;\le\;
\frac{2\bigl(f(x_{0})-f^{\star}\bigr)}{\eta T}
\;+\;
L\eta\bigl(C^{2}+\zeta^{2}+d\sigma^{2}\bigr),
\tag{2}
\]
where $f^{\star}=\inf_{x}f(x)$.
Consequently, choosing $\eta=\Theta\!\bigl(T^{-1/2}\bigr)$ yields
\[
\frac{1}{T}\sum_{t=0}^{T-1}\E\!\bigl[\|\grad f(x_{t})\|^{2}\bigr]=\mathcal{O}\!\bigl(T^{-1/2}\bigr).
\]
\end{theorem}

\section*{Theoretical Properties}
\begin{itemize}
\item \textbf{Stability.} The clipping bound $C$ guarantees that the variance contributed by the data term is uniformly bounded, preventing exploding updates even for non‑convex losses.
\item \textbf{Hyperparameter Sensitivity.} The bound (2) shows a linear dependence on $\eta$ for the variance term (which contains $C^{2},\zeta^{2},d\sigma^{2}$) and an inverse dependence on $\eta$ for the optimization term. The optimal $\eta$ balances these two effects, leading to the $\mathcal{O}(T^{-1/2})$ rate.
\item \textbf{Comparison to Baselines.} For $\sigma=0$ (no privacy) the theorem reduces to the classical non‑convex SGD bound $\mathcal{O}(T^{-1/2})$ \cite{Ghadimi2013}. Adding Gaussian noise inflates the variance term by $L\eta d\sigma^{2}$, matching known DP‑SGD analyses.
\item \textbf{Privacy–Utility Trade‑off.} Larger $\sigma$ (stronger privacy) linearly worsens the stationary‑point guarantee, while a smaller $C$ reduces the data‑dependent variance at the cost of potentially biasing the gradient direction.
\end{itemize}

\section*{Discussion}
The theorem guarantees that DP‑SGD finds an \emph{approximate stationary point} of a non‑convex smooth objective at the same $\mathcal{O}(T^{-1/2})$ rate as vanilla SGD, up to an additive term proportional to the privacy noise variance. This aligns with the intuition that differential privacy introduces additional stochasticity but does not change the fundamental convergence order. The proof below follows the standard smoothness‑based descent argument, carefully tracking the extra variance from clipping and Gaussian noise.

\section*{Preliminaries (Definitions and Lemmas)}
\begin{lemma}[Smoothness Inequality]\label{lem:smooth}
If $f$ is $L$‑smooth then for any $x,y\in\R^{d}$,
\[
f(y)\le f(x)+\inner{\grad f(x)}{y-x}+\frac{L}{2}\|y-x\|^{2}.
\]
\end{lemma}
\begin{lemma}[Variance Decomposition]\label{lem:var}
For any random vector $u$ with $\E[u]=\mu$, 
\[
\E[\|u\|^{2}]=\|\mu\|^{2}+\E[\|u-\mu\|^{2}].
\]
\end{lemma}

\section*{Main Proof (Step‑by‑Step Derivation)}
We prove Theorem \ref{thm:main}. All expectations are taken over the randomness of the mini‑batch, the clipping, and the added Gaussian noise unless otherwise noted.

\noindent\textbf{[Step 1]} Apply Lemma \ref{lem:smooth} with $x=x_{t}$ and $y=x_{t+1}=x_{t}-\eta g_{t}$:
\[
f(x_{t+1})\le f(x_{t})-\eta\inner{\grad f(x_{t})}{g_{t}}+\frac{L\eta^{2}}{2}\|g_{t}\|^{2}.
\tag{3}
\]

\noindent\textbf{[Step 2]} Expand $g_{t}=\bar g_{t}+z_{t}$ and use linearity of inner product:
\[
\inner{\grad f(x_{t})}{g_{t}}
= \inner{\grad f(x_{t})}{\bar g_{t}} + \inner{\grad f(x_{t})}{z_{t}}.
\tag{4}
\]

\noindent\textbf{[Step 3]} Condition on $x_{t}$. By Assumption \ref{as:unbiased}, $\E[\bar g_{t}\mid x_{t}]=\grad f(x_{t})$, and by Assumption \ref{as:noise}, $\E[z_{t}]=0$. Hence
\[
\E\!\bigl[\inner{\grad f(x_{t})}{\bar g_{t}} \mid x_{t}\bigr]=\|\grad f(x_{t})\|^{2},
\qquad
\E\!\bigl[\inner{\grad f(x_{t})}{z_{t}}\mid x_{t}\bigr]=0.
\tag{5}
\]

\noindent\textbf{[Step 4]} Take total expectation of (3) and use (5):
\[
\E\!\bigl[f(x_{t+1})\bigr]
\le \E\!\bigl[f(x_{t})\bigr]
-\eta\E\!\bigl[\|\grad f(x_{t})\|^{2}\bigr]
+\frac{L\eta^{2}}{2}\E\!\bigl[\|g_{t}\|^{2}\bigr].
\tag{6}
\]

\noindent\textbf{[Step 5]} Decompose $\|g_{t}\|^{2}$ via Lemma \ref{lem:var} with $u=g_{t}$, $\mu=\E[g_{t}\mid x_{t}]=\grad f(x_{t})$ (by (5)):
\[
\E\!\bigl[\|g_{t}\|^{2}\mid x_{t}\bigr]
= \|\grad f(x_{t})\|^{2}
+ \E\!\bigl[\|g_{t}-\grad f(x_{t})\|^{2}\mid x_{t}\bigr].
\tag{7}
\]

\noindent\textbf{[Step 6]} Write $g_{t}-\grad f(x_{t})=(\bar g_{t}-\grad f(x_{t}))+z_{t}$. Using independence of $z_{t}$ from the mini‑batch,
\[
\E\!\bigl[\|g_{t}-\grad f(x_{t})\|^{2}\mid x_{t}\bigr]
= \E\!\bigl[\|\bar g_{t}-\grad f(x_{t})\|^{2}\mid x_{t}\bigr]
+ \E\!\bigl[\|z_{t}\|^{2}\bigr].
\tag{8}
\]

\noindent\textbf{[Step 7]} Bound the variance of the clipped mini‑batch gradient. By definition of $\bar g_{t}$,
\[
\bar g_{t}= \frac{1}{b}\sum_{i\in\mathcal{B}_{t}}\tilde g_{i}^{t},
\qquad
\tilde g_{i}^{t}= \operatorname{clip}(\grad f_{i}(x_{t}),C).
\]
Since each $\tilde g_{i}^{t}$ has norm at most $C$, the variance of their average satisfies
\[
\E\!\bigl[\|\bar g_{t}-\grad f(x_{t})\|^{2}\mid x_{t}\bigr]
\le \frac{1}{b}\bigl(C^{2}+\zeta^{2}\bigr),
\tag{9}
\]
where the term $C^{2}$ comes from the boundedness of each clipped gradient and $\zeta^{2}$ from Assumption \ref{as:bounded} (the second‑moment deviation of a single example).  
(For a rigorous derivation one expands the variance of an average of i.i.d. bounded random vectors.)

\noindent\textbf{[Step 8]} The Gaussian noise variance is $\E[\|z_{t}\|^{2}]=d\sigma^{2}$ by Assumption \ref{as:noise}. Plugging (9) and this into (8) yields
\[
\E\!\bigl[\|g_{t}-\grad f(x_{t})\|^{2}\mid x_{t}\bigr]
\le \frac{1}{b}\bigl(C^{2}+\zeta^{2}\bigr)+d\sigma^{2}.
\tag{10}
\]

\noindent\textbf{[Step 9]} Substitute (10) into (7) and then into (6):
\[
\E\!\bigl[f(x_{t+1})\bigr]
\le \E\!\bigl[f(x_{t})\bigr]
-\eta\E\!\bigl[\|\grad f(x_{t})\|^{2}\bigr]
+\frac{L\eta^{2}}{2}\Bigl(\E\!\bigl[\|\grad f(x_{t})\|^{2}\bigr]
+\frac{1}{b}\bigl(C^{2}+\zeta^{2}\bigr)+d\sigma^{2}\Bigr).
\tag{11}
\]

\noindent\textbf{[Step 10]} Rearrange (11) to isolate the gradient‑norm term:
\[
\Bigl(\eta-\frac{L\eta^{2}}{2}\Bigr)\E\!\bigl[\|\grad f(x_{t})\|^{2}\bigr]
\le \E\!\bigl[f(x_{t})-f(x_{t+1})\bigr]
+\frac{L\eta^{2}}{2}\Bigl(\frac{1}{b}\bigl(C^{2}+\zeta^{2}\bigr)+d\sigma^{2}\Bigr).
\tag{12}
\]

\noindent\textbf{[Step 11]} Since $\eta\le 1/L$ (Assumption \ref{as:lr}), we have $\eta-\frac{L\eta^{2}}{2}\ge \frac{\eta}{2}$. Hence
\[
\frac{\eta}{2}\E\!\bigl[\|\grad f(x_{t})\|^{2}\bigr]
\le \E\!\bigl[f(x_{t})-f(x_{t+1})\bigr]
+\frac{L\eta^{2}}{2}\Bigl(\frac{1}{b}\bigl(C^{2}+\zeta^{2}\bigr)+d\sigma^{2}\Bigr).
\tag{13}
\]

\noindent\textbf{[Step 12]} Sum inequality (13) over $t=0,\dots,T-1$ and telescope the objective terms:
\[
\frac{\eta}{2}\sum_{t=0}^{T-1}\E\!\bigl[\|\grad f(x_{t})\|^{2}\bigr]
\le f(x_{0})- \E\!\bigl[f(x_{T})\bigr]
+ \frac{L\eta^{2}T}{2}\Bigl(\frac{1}{b}\bigl(C^{2}+\zeta^{2}\bigr)+d\sigma^{2}\Bigr).
\tag{14}
\]

\noindent\textbf{[Step 13]} Drop the non‑negative term $-\E[f(x_{T})]\le 0$ and divide by $\eta T$:
\[
\frac{1}{T}\sum_{t=0}^{T-1}\E\!\bigl[\|\grad f(x_{t})\|^{2}\bigr]
\le \frac{2\bigl(f(x_{0})-f^{\star}\bigr)}{\eta T}
+ L\eta\Bigl(\frac{1}{b}\bigl(C^{2}+\zeta^{2}\bigr)+d\sigma^{2}\Bigr).
\tag{15}
\]

\noindent\textbf{[Step 14]} Recognize that $f^{\star}\le \E[f(x_{T})]$, yielding the bound stated in Theorem \ref{thm:main}. This completes the proof. \hfill $\square$

\section*{Rate Derivation (With Constants)}
Choosing $\eta = \sqrt{\frac{2\bigl(f(x_{0})-f^{\star}\bigr)}{L\bigl(C^{2}+\zeta^{2}+bd\sigma^{2}\bigr)T}}$ (which satisfies $\eta\le 1/L$ for sufficiently large $T$) balances the two terms in (15) and gives
\[
\frac{1}{T}\sum_{t=0}^{T-1}\E\!\bigl[\|\grad f(x_{t})\|^{2}\bigr]
\le
\underbrace{2\sqrt{\frac{L\bigl(C^{2}+\zeta^{2}+bd\sigma^{2}\bigr)
\bigl(f(x_{0})-f^{\star}\bigr)}{T}}}_{\displaystyle \mathcal{O}(T^{-1/2})}.
\]
Thus the algorithm attains the canonical $\mathcal{O}(T^{-1/2})$ stationary‑point rate, with a constant factor that explicitly captures the effect of clipping $C$, data variance $\zeta^{2}$, batch size $b$, dimension $d$, and privacy noise $\sigma^{2}$.

\section*{Special Cases}
\begin{itemize}
\item \textbf{No privacy noise ($\sigma=0$).} The bound reduces to the standard non‑convex SGD result with variance term $\frac{L\eta}{b}(C^{2}+\zeta^{2})$.
\item \textbf{Full‑batch ($b=n$) and no clipping ($C=\infty$).} Then $C^{2}$ term disappears, $\zeta^{2}=0$, and (15) becomes the deterministic gradient descent bound $\frac{2(f(x_{0})-f^{\star})}{\eta T}+L\eta d\sigma^{2}$.
\item \textbf{High privacy ($\sigma$ large).} The term $L\eta d\sigma^{2}$ dominates; the optimal step size scales as $\eta\propto T^{-1/2}\sigma^{-1}$, reflecting the trade‑off between privacy and convergence speed.
\end{itemize}

\begin{thebibliography}{9}
\bibitem{Ghadimi2013}
S. Ghadimi and G. Lan, ``Stochastic First‑and Zero‑Order Methods for Nonconvex Stochastic Programming,'' \emph{SIAM J. Optim.}, vol. 23, no. 4, pp. 2341–2368, 2013.
\end{thebibliography}

\end{document}